\makeatletter
\def\input@path{{./}{../}{../../}{../../../}{../../../../}{preamble/}{../preamble/}{../../preamble/}{../../../preamble/}}
\makeatother

\input{preamble_exams.tex}
\input{graphs/highlights.tex}
% \GTenHighlightsfalse % Uncomment to disable instructional highlights.
\input{graphs/specific_highlights.tex}
\CoordinatedTermTwoPalette
\usepackage{needspace}
\usepackage{titling}
\setlength{\droptitle}{-2em}
\pretitle{\par\begingroup\centering\LARGE}
\posttitle{\par\endgroup}
\preauthor{}\postauthor{}
\predate{}\postdate{}
\maketitlehookd{\vspace{-1em}}
\linespread{1.12}
\setlength{\parskip}{5pt}
\newgeometry{margin=0.65in,top=0.55in,bottom=0.65in}

\title{Grade 11 Term 2 Home Practice -- C6\\\large Planning sample size\\[0.25em]
	\normalsize Global Economics -- 2026--2027}
\author{}
\date{}


% Local specialization of the existing C6 entry: shared registry stays unchanged.
\expandafter\renewcommand\csname practicehighlightlegend@G11-T2-C6-coordinated\endcsname{%
\begingroup
\CoordinatedLegendTable{
$\HLObservationCount{n}$, $\HLObservationCount{n_{\min}}$ & \HLObservationCount{Sample size} & Number of usable observations.\\
$\HLPopulationSD{\sigma}$; $\HLModelControl{p_0}$ & Fixed variability inputs & Known spread; assumed planning share.\\
$\HLStandardError{SE(n)}$, $\HLStandardError{S}$ & \HLStandardError{Standard error} & Sampling spread and its cap.\\
$\HLConfidence{C}$, $\HLConfidence{C_0}$ & \HLConfidence{Confidence level} & Central normal probability and target.\\
$\HLCriticalValue{z^*}$, $\HLCriticalValue{z_0}$ & \HLCriticalValue{Critical value} & Normal multiplier for confidence.\\
$\HLMarginOfError{E}$, $\HLMarginOfError{E_0}$ & \HLMarginOfError{Margin of error} & Half-width and its cap.\\
$\HLIntervalWidth{W}$, $\HLIntervalWidth{W_0}$ & \HLIntervalWidth{Total width} & Full interval length and its cap.\\
}\endgroup}
\newcommand{\IntroOne}{\textbullet\ \hyperref[ex:intro-one]{Example 1}: }
\newcommand{\IntroTwo}{\textbullet\ \hyperref[ex:intro-two]{Example 2}: }

\begin{document}
\maketitle
\setcounter{tocdepth}{1}\setcounter{secnumdepth}{3}
\phantomsection\label{toc}\tableofcontents
\section*{Evaluated criteria}
\begin{enumerate}[label=\textbf{C\arabic*},start=6]
\item Determines the sample size required to estimate a population's mean or proportion with a specified margin of error, using it to plan economic surveys or financial studies with predetermined precision.
\end{enumerate}
\textbf{How to practise.} Try each task before opening its blue solution link. Record the target, what stays fixed, the inequality, the unrounded bound and the smallest integer that works. The arrows return to the contents, problem list or current problem.
\clearpage
\section{Introduction}\label{sec:introduction}
\subsection*{Highlight legend}\phantomsection\label{tab:highlight-legend}
\PracticeHighlightLegend{G11-T2-C6-coordinated}
Names and units retain meaning without color. Related shades connect related quantities; standard error, confidence, critical value, half-width and total width remain distinct roles.
\subsection{Introductory example 1: cargo clearance}\label{ex:intro-one}
A port wants the unknown population mean $\HLPopulationMean{\mu}$ clearance time. Its known population standard deviation is $\HLPopulationSD{\sigma}=\HLPopulationSD{12\text{ min}}$. For a random sample of $\HLObservationCount{n}$ records, the sample mean has
\[
\HLStandardError{SE_{\rm mean}(n)}=\frac{\HLPopulationSD{\sigma}}{\sqrt{\HLObservationCount{n}}}.
\]
At $n=100$, $SE=1.2$ minutes. A central $\HLConfidence{95\%}$ normal interval uses $\HLCriticalValue{z^*}\approx\HLCriticalValue{1.95996398454}$, so its $\HLMarginOfError{E}=z^*SE\approx2.352$ minutes and its $\HLIntervalWidth{W}=2E\approx4.704$ minutes. The interval is centered on the future sample mean, not on the unknown $\mu$.
\subsection{Introductory example 2: household energy}\label{ex:intro-two}
An energy office wants the unknown population proportion $\HLPopulationProportion{p}$ eligible for a subsidy. A pilot suggests the fixed planning value $\HLModelControl{p_0}=\HLModelControl{0.20}$; the future sample proportion $\HLSampleProportion{\hat p}$ will be calculated from collected responses. For planning,
\[
\HLStandardError{SE_{\rm proportion}(n)}\approx
\sqrt{\frac{\HLModelControl{p_0}(1-\HLModelControl{p_0})}{\HLObservationCount{n}}}.
\]
At $n=400$, planning $SE=0.02$, or 2 percentage points. A 95\% interval has planning $E\approx0.03920$ (3.920 percentage points), and $W\approx0.07840$ (7.840 percentage points). A cap of 4 percentage points means $E_0=0.04$, not 4. The pilot is an assumption, not knowledge of $p$ or a prediction that $\hat p=p_0$.
\textbf{Quick check.} A full width of 8 percentage points permits a margin of only 4 percentage points. Changing confidence changes the multiplier, not the standard error for a fixed design and $n$.
\clearpage
\section{Notation Reference and Practice Procedure}\label{sec:notation-and-procedure}
\subsection{Notation reference}\label{tab:notation-reference}
\begingroup\tikzset{gten marker/.append style={inner xsep=0.6mm,inner ysep=0.3mm}}\setlength{\tabcolsep}{4pt}\renewcommand{\arraystretch}{1.12}
\begin{longtable}{@{}>{\raggedright\arraybackslash}p{.17\linewidth}>{\raggedright\arraybackslash}p{.18\linewidth}>{\raggedright\arraybackslash}p{.27\linewidth}>{\raggedright\arraybackslash}p{\dimexpr.38\linewidth-6\tabcolsep\relax}@{}}
\toprule
\textbf{Symbol or notation}&\textbf{Concept}&\textbf{Definition}&\textbf{Meaning in the introductory examples}\\\midrule\endfirsthead
\toprule\textbf{Symbol or notation}&\textbf{Concept}&\textbf{Definition}&\textbf{Meaning in the introductory examples}\\\midrule\endhead
$\HLObservationCount{n},\HLObservationCount{n_{\min}}$ & \HLObservationCount{Sample size} & Positive integer count; least count meeting all targets. & \IntroOne Usable cargo records.\par\IntroTwo Completed household responses.\\
$\HLPopulationMean{\mu},\HLSampleMean{\bar x}$ & Population and sample means & Fixed population average; average calculated from the sample. & \IntroOne Unknown mean clearance time and its estimate, in minutes.\\
$\HLPopulationSD{\sigma}$ & \HLPopulationSD{Population} \HLPopulationSD{spread} & Known population standard deviation, in response units. & \IntroOne $\HLPopulationSD{12\text{ min}}$, held fixed as $n$ changes.\\
$\HLPopulationProportion{p}$; $\HLModelControl{p_0}$; $\HLSampleProportion{\hat p}$ & Population, planning and sample shares & Unknown population share; assumed design input; observed sample share. All are proportions. & \IntroTwo Subsidy eligibility: $p$ is unknown, $\HLModelControl{p_0}=\HLModelControl{0.20}$ comes from a pilot; $\hat p$ is future data.\\
$\HLStandardError{SE(n)},\HLStandardError{S}$ & \HLStandardError{Sampling} \HLStandardError{spread} and cap & Standard deviation of an estimator (planning approximation for a proportion); maximum allowed SE. & \IntroOne Minutes.\par\IntroTwo Proportion units; $0.02=2$ percentage points.\\
$\HLConfidence{C},\HLConfidence{C_0}$ & \HLConfidence{Confidence} and target & Central normal probability and its minimum required level. & \IntroOne $\HLConfidence{95\%}$ confidence.\par\IntroTwo $\HLConfidence{0.95}$ as a proportion.\\
$\HLCriticalValue{z^*},\HLCriticalValue{z_0}$; $\Phi$ & \HLCriticalValue{Normal} \HLCriticalValue{quantile}; CDF & $z^*=\Phi^{-1}((1+C)/2)$; $z_0$ uses $C_0$. $\Phi(t)=P(Z\le t)$ for $Z\sim N(0,1)$. & \IntroOne The 95\% multiplier.\par\IntroTwo The same multiplier, $\HLCriticalValue{z^*}\approx\HLCriticalValue{1.95996398454}$.\\
$\HLMarginOfError{E},\HLMarginOfError{E_0}$ & \HLMarginOfError{Half-width} and cap & $E=z^*SE$; maximum permitted distance from interval center to an endpoint. & \IntroOne Minutes.\par\IntroTwo $\HLMarginOfError{0.04}$ means a cap of $\HLMarginOfError{4}$ percentage points.\\
$\HLIntervalWidth{W},\HLIntervalWidth{W_0}$ & \HLIntervalWidth{Total width} and cap & $W=2E$; maximum full interval length. & \IntroOne About 4.704 minutes.\par\IntroTwo About 7.840 percentage points.\\
$V$; $\lceil b\rceil$ & Compact variability; ceiling & $V=\sigma^2$ or $p_0(1-p_0)$; least integer at least $b$. & \IntroOne $V=144\text{ min}^2$.\par\IntroTwo $V=0.16$ (squared proportion units).\\
$N$; $np_0$, $n(1-p_0)$ & Population size; expected counts & Use $n\le0.10N$ without replacement; check both expected counts at least 10. & \IntroTwo Expected eligible and other households under the planning assumption.\\
\bottomrule\end{longtable}\endgroup

\clearpage
\subsection{Models, assumptions and calculator use}
For a mean, use a simple random sample, known $\sigma>0$, and a normal population or a sample mean for which a normal approximation is suitable. For a proportion, use random binary outcomes and the normal planning approximation, with $0<p_0<1$. A pilot-based result is conditional on that assumed planning proportion. Check $np_0\ge10$ and $n(1-p_0)\ge10$ at planning; after collection check the actual counts $n\hat p$ and $n(1-\hat p)$ and the interval conditions again.

Without a dependable pilot, $p_0=0.50$ is conservative because
\[
p_0(1-p_0)=\tfrac14-(p_0-\tfrac12)^2\le\tfrac14.
\]
This bounds variability and the normal-formula planning margin over possible proportions. It does not guarantee exact finite-sample coverage or accuracy for every realized sample. All worked populations below have $N=100{,}000$ eligible units; their samples satisfy $n\le10{,}000$. This supports approximate independence without replacement. Finite-population corrections are outside scope.

\textbf{Calculator instructions.} Use inverse normal with mean 0, standard deviation 1 and left-tail area $(1+C)/2$. Use normal CDF with lower endpoint $-\infty$, upper endpoint $t$, mean 0 and standard deviation 1 to obtain $\Phi(t)$. Store full calculator precision:
\[
z_{.95}=\Phi^{-1}(0.975)\approx1.959963984540054,\qquad
z_{.90}=\Phi^{-1}(0.95)\approx1.644853626951472.
\]
These stored values are used throughout; displayed results are rounded only after calculation. Use extra digits when comparing a result with a cap; a coarse normal table can obscure a minimum-size decision.
\subsection{Reusable decision procedure}
\begin{enumerate}
\item Identify the unknown integer $n$, the estimator and units. List the target and what stays fixed: variability always; confidence in Types 2--3; precision cap in Type 4. Convert percentage points to proportions.
\item Write the explicit mean or proportion SE formula first. Check the model, sampling fraction and expected counts where relevant.
\item Choose the inequality: $SE\le S$, $z^*SE\le E_0$, $2z^*SE\le W_0$, or $C(n)\ge C_0$ with fixed $E_0>0$.
\item Rearrange using positive quantities. Calculate the unrounded lower bound $b$, then $n_{\min}=\lceil b\rceil$. An integer bound stays unchanged; never automatically add one.
\item Check $n_{\min}$ works and $n_{\min}-1$ fails when $n_{\min}>1$. The relevant SE, margin and width decrease with $n$; fixed-cap confidence increases with $n$. Thus the adjacent checks establish minimality.
\item Interpret the answer as usable observations or completed responses. Invitations may need to exceed this count if there is nonresponse.
\end{enumerate}
\clearpage
\section{General Worked Problems}\label{sec:worked-solution-examples}
Choose a question type: \hyperref[prob:t1]{standard error},
\hyperref[prob:t2]{margin of error}, \hyperref[prob:t3]{total width}, or
\hyperref[prob:t4]{confidence with fixed precision}.

After the explicit SE formulas above, write $V=\sigma^2$ for a mean or $V=p_0(1-p_0)$ for a proportion. Thus $SE(n)=\sqrt{V/n}$ (approximately for a proportion). Treat all proportion calculations below as normal-formula planning calculations.

\subsection{Type 1: minimum sample size for a standard-error cap}\label{prob:t1}
Keep population variability fixed and supply $\HLStandardError{S}>0$. Confidence level does not enter either standard-error formula:
\[
\sqrt{V/n}\le S\ \Longleftrightarrow\ V/n\le S^2
\ \Longleftrightarrow\ n\ge V/S^2,
\qquad n_{\min}=\left\lceil V/S^2\right\rceil.
\]
\SectionProblemsTableOfContents{
\SectionProblemLink{prob:t1-mean}{Mean: cargo records}\\
\SectionProblemLink{prob:t1-prop}{Proportion: pilot-based planning}}
\medskip
\subsubsection{Cargo records: a standard error of at most 1.5 minutes}\label{prob:t1-mean}
A logistics agency estimates the unknown population mean cargo-clearance time. It will draw a simple random sample without replacement from $N=100{,}000$ eligible cargo records. Hold the known population spread $\HLPopulationSD{\sigma}=\HLPopulationSD{12\text{ min}}$ fixed. Assume the sample mean is suitably approximated by a normal distribution. The standard error cap is 1.5 minutes. No confidence level is needed for this standard-error target. 
\QuestionsTasks[sec:worked-solution-examples]{{6}/{sol:t1-mean}/{Find the minimum integer sample size. State the unknown, target, fixed quantities, units and model; rearrange the inequality, round correctly, check adjacent integers and interpret the result.}}
\subsubsection*{Worked solution (C6)}\label{sol:t1-mean}
\textbf{Design.} The unknown is $\HLObservationCount{n}$, the number of usable cargo records. The target is standard error at most 1.5 minutes. Population variability remains fixed. Use the known-$\sigma$ normal model for the sample mean. Confidence does not enter $SE(n)$. 
\[
\frac{\HLPopulationSD{12}}{\sqrt n}\le \HLStandardError{1.5}\quad\Longrightarrow\quad\sqrt n\ge\frac{\HLPopulationSD{12}}{1.5}.
\]
\[n\ge\left(\frac{\HLPopulationSD{12}}{1.5}\right)^2=64,\qquad \HLObservationCount{n_{\min}}=\left\lceil 64\right\rceil=\HLObservationCount{64}.\]
The bound is already an integer, so retain it; do not add one.
\textbf{Adjacent-integer check} (minutes):
\[ SE(64)\approx1.500000000\le \HLStandardError{1.5},\qquad SE(63)\approx1.511857892>1.5.\]
The sampling fraction is $64/100{,}000=0.064\%<10\%$. Because the target quantity decreases with $n$, the checks establish minimality. Plan for $\HLObservationCount{64}$ usable cargo records. 
\secInicio[prob:t1-mean][sec:worked-solution-examples]
\clearpage
\subsubsection{Energy eligibility: a standard error of at most 2.5 percentage points}\label{prob:t1-prop}
An energy office estimates the unknown population share $p$ of households eligible for a bill subsidy. It will draw a simple random sample without replacement from $N=100{,}000$ households. A pilot suggests $\HLModelControl{p_0}=\HLModelControl{0.20}$; hold this planning value fixed. The standard error cap is 2.5 percentage points, namely 0.025 in proportion units. No confidence level is needed for this standard-error target. 
\QuestionsTasks[sec:worked-solution-examples]{{6}/{sol:t1-prop}/{Find the minimum integer sample size. State the unknown, target, fixed quantities, units and model; rearrange the inequality, round correctly, check adjacent integers and interpret the result.}}
\subsubsection*{Worked solution (C6)}\label{sol:t1-prop}
\textbf{Design.} The unknown is $\HLObservationCount{n}$, the number of usable household responses. The target is standard error at most 2.5 percentage points, namely 0.025 in proportion units. The planning share remains fixed; $p$ is unknown and $\hat p$ is the future estimate. Use the normal approximation for random binary outcomes; expected-count checks appear below. Confidence does not enter $SE(n)$. 
\[
\sqrt{\frac{\HLModelControl{0.20}(0.80)}{n}}\le \HLStandardError{0.025}\quad\Longrightarrow\quad\sqrt n\ge\frac{\sqrt{\HLModelControl{0.20}(0.80)}}{0.025}.
\]
\[n\ge\left(\frac{\sqrt{\HLModelControl{0.20}(0.80)}}{0.025}\right)^2=256,\qquad \HLObservationCount{n_{\min}}=\left\lceil 256\right\rceil=\HLObservationCount{256}.\]
The bound is already an integer, so retain it; do not add one.
\textbf{Adjacent-integer check} (proportion units):
\[ SE(256)\approx0.025000000\le \HLStandardError{0.025},\qquad SE(255)\approx0.025048972>0.025.\]
Expected outcome counts are $np_0=51.2$ and $n(1-p_0)=204.8$, both at least 10. The sampling fraction is $256/100{,}000=0.256\%<10\%$. Because the target quantity decreases with $n$, the checks establish minimality. Plan for $\HLObservationCount{256}$ usable household responses. This pilot-based result is conditional on the assumed planning proportion. After data collection, check actual outcome counts and the interval conditions. 
\secInicio[prob:t1-prop][sec:worked-solution-examples]

\clearpage
\subsection{Type 2: minimum sample size for a margin-of-error cap}\label{prob:t2}
Keep confidence level and population variability fixed; supply $\HLMarginOfError{E_0}>0$:
\[
z^*\sqrt{V/n}\le E_0
\ \Longleftrightarrow\ \sqrt n\ge\frac{z^*\sqrt V}{E_0}
\ \Longleftrightarrow\ n\ge\frac{(z^*)^2V}{E_0^2}.
\]
Hence $n_{\min}=\left\lceil (z^*)^2V/E_0^2\right\rceil$.
\SectionProblemsTableOfContents{
\SectionProblemLink{prob:c6-1}{Mean: cargo records}\\
\SectionProblemLink{prob:c6-2}{Proportion: pilot-based planning}}
\medskip
\subsubsection{Planning a port inspection study (C6)}\label{prob:c6-1}
A logistics agency estimates the unknown population mean cargo-clearance time. It will draw a simple random sample without replacement from $N=100{,}000$ eligible cargo records. Hold the known population spread $\HLPopulationSD{\sigma}=\HLPopulationSD{12\text{ min}}$ fixed. Assume the sample mean is suitably approximated by a normal distribution. The margin of error cap is 3 minutes. Hold confidence fixed at $\HLConfidence{95\%}$, using $z^*=z_{.95}$. 
\QuestionsTasks[sec:worked-solution-examples]{{6}/{sol:c6-1}/{Find the minimum integer sample size. State the unknown, target, fixed quantities, units and model; rearrange the inequality, round correctly, check adjacent integers and interpret the result.}}
\subsubsection*{Worked solution (C6)}\label{sol:c6-1}
\phantomsection\label{sol:c6-1-a}\label{sol:c6-1-b}
\textbf{Design.} The unknown is $\HLObservationCount{n}$, the number of usable cargo records. The target is margin of error at most 3 minutes with 95\% confidence. Population variability remains fixed. Use the known-$\sigma$ normal model for the sample mean. Confidence also remains fixed. 
\[
\HLCriticalValue{z^*}\frac{\HLPopulationSD{12}}{\sqrt n}\le \HLMarginOfError{3}\quad\Longrightarrow\quad\sqrt n\ge\frac{\HLCriticalValue{z^*}\HLPopulationSD{12}}{3}.
\]
\[n\ge\left(\frac{\HLCriticalValue{z^*}\HLPopulationSD{12}}{3}\right)^2\approx61.463341131,\qquad \HLObservationCount{n_{\min}}=\left\lceil 61.463341131\right\rceil=\HLObservationCount{62}.\]
Round upward, because rounding down would violate the inequality.
\textbf{Adjacent-integer check} (minutes):
\[ E(62)\approx2.986988099\le \HLMarginOfError{3},\qquad E(61)\approx3.011372080>3.\]
The sampling fraction is $62/100{,}000=0.062\%<10\%$. Because the target quantity decreases with $n$, the checks establish minimality. Plan for $\HLObservationCount{62}$ usable cargo records. 
\secInicio[prob:c6-1][sec:worked-solution-examples]
\clearpage
\subsubsection{Planning a household energy survey (C6)}\label{prob:c6-2}
An energy office estimates the unknown population share $p$ of households eligible for a bill subsidy. It will draw a simple random sample without replacement from $N=100{,}000$ households. A pilot suggests $\HLModelControl{p_0}=\HLModelControl{0.20}$; hold this planning value fixed. The margin of error cap is 4 percentage points, namely 0.04 in proportion units. Hold confidence fixed at $\HLConfidence{95\%}$, using $z^*=z_{.95}$. 
\QuestionsTasks[sec:worked-solution-examples]{{6}/{sol:c6-2}/{Find the minimum integer sample size. State the unknown, target, fixed quantities, units and model; rearrange the inequality, round correctly, check adjacent integers and interpret the result.}}
\subsubsection*{Worked solution (C6)}\label{sol:c6-2}
\phantomsection\label{sol:c6-2-a}\label{sol:c6-2-b}
\textbf{Design.} The unknown is $\HLObservationCount{n}$, the number of usable household responses. The target is margin of error at most 4 percentage points, namely 0.04 in proportion units with 95\% confidence. The planning share remains fixed; $p$ is unknown and $\hat p$ is the future estimate. Use the normal approximation for random binary outcomes; expected-count checks appear below. Confidence also remains fixed. 
\[
\HLCriticalValue{z^*}\sqrt{\frac{\HLModelControl{0.20}(0.80)}{n}}\le \HLMarginOfError{0.04}\quad\Longrightarrow\quad\sqrt n\ge\frac{\HLCriticalValue{z^*}\sqrt{\HLModelControl{0.20}(0.80)}}{0.04}.
\]
\[n\ge\left(\frac{\HLCriticalValue{z^*}\sqrt{\HLModelControl{0.20}(0.80)}}{0.04}\right)^2\approx384.145882069,\qquad \HLObservationCount{n_{\min}}=\left\lceil 384.145882069\right\rceil=\HLObservationCount{385}.\]
Round upward, because rounding down would violate the inequality.
\textbf{Adjacent-integer check} (proportion units):
\[ E(385)\approx0.039955606\le \HLMarginOfError{0.04},\qquad E(384)\approx0.040007597>0.04.\]
Expected outcome counts are $np_0=77$ and $n(1-p_0)=308$, both at least 10. The sampling fraction is $385/100{,}000=0.385\%<10\%$. Because the target quantity decreases with $n$, the checks establish minimality. Plan for $\HLObservationCount{385}$ usable household responses. This pilot-based result is conditional on the assumed planning proportion. After data collection, check actual outcome counts and the interval conditions. 
\secInicio[prob:c6-2][sec:worked-solution-examples]

\clearpage
\subsection{Type 3: minimum sample size for a total-width cap}\label{prob:t3}
Keep confidence level and population variability fixed; supply $\HLIntervalWidth{W_0}>0$. Total width is twice the margin, so $E_0=W_0/2$:
\[
2z^*\sqrt{V/n}\le W_0
\ \Longleftrightarrow\ \sqrt n\ge\frac{2z^*\sqrt V}{W_0}
\ \Longleftrightarrow\ n\ge\frac{4(z^*)^2V}{W_0^2}.
\]
Hence $n_{\min}=\left\lceil4(z^*)^2V/W_0^2\right\rceil$.
\SectionProblemsTableOfContents{
\SectionProblemLink{prob:t3-mean}{Mean: export records}\\
\SectionProblemLink{prob:t3-prop}{Proportion: conservative planning}}
\medskip
\subsubsection{Export processing: a total width of at most 5 minutes}\label{prob:t3-mean}
A logistics agency estimates the unknown population mean export-processing time. It will draw a simple random sample without replacement from $N=100{,}000$ eligible export records. Hold the known population spread $\HLPopulationSD{\sigma}=\HLPopulationSD{15\text{ min}}$ fixed. Assume the sample mean is suitably approximated by a normal distribution. The total interval width cap is 5 minutes. Hold confidence fixed at $\HLConfidence{90\%}$, using $z^*=z_{.90}$. 
\QuestionsTasks[sec:worked-solution-examples]{{6}/{sol:t3-mean}/{Find the minimum integer sample size. State the unknown, target, fixed quantities, units and model; rearrange the inequality, round correctly, check adjacent integers and interpret the result.}}
\subsubsection*{Worked solution (C6)}\label{sol:t3-mean}
\textbf{Design.} The unknown is $\HLObservationCount{n}$, the number of usable export records. The target is total interval width at most 5 minutes with 90\% confidence. Population variability remains fixed. Use the known-$\sigma$ normal model for the sample mean. Confidence also remains fixed. 
\[
2\HLCriticalValue{z^*}\frac{\HLPopulationSD{15}}{\sqrt n}\le \HLIntervalWidth{5}\quad\Longrightarrow\quad\sqrt n\ge\frac{2\HLCriticalValue{z^*}\HLPopulationSD{15}}{5}.
\]
\[n\ge\left(\frac{2\HLCriticalValue{z^*}\HLPopulationSD{15}}{5}\right)^2\approx97.399564347,\qquad \HLObservationCount{n_{\min}}=\left\lceil 97.399564347\right\rceil=\HLObservationCount{98}.\]
Round upward, because rounding down would violate the inequality.
\textbf{Adjacent-integer check} (minutes):
\[ W(98)\approx4.984659230\le \HLIntervalWidth{5},\qquad W(97)\approx5.010287467>5.\]
The permitted half-width is $2.5$ minutes; $5$ is the full width, not the margin. The sampling fraction is $98/100{,}000=0.098\%<10\%$. Because the target quantity decreases with $n$, the checks establish minimality. Plan for $\HLObservationCount{98}$ usable export records. 
\secInicio[prob:t3-mean][sec:worked-solution-examples]
\clearpage
\subsubsection{Energy uptake: a total width of at most 6 percentage points}\label{prob:t3-prop}
An energy office estimates the unknown population share $p$ of households adopting an efficiency tariff. It will draw a simple random sample without replacement from $N=100{,}000$ households. No dependable pilot is available: hold $\HLModelControl{p_0}=\HLModelControl{0.50}$ fixed for conservative planning. The total interval width cap is 6 percentage points, namely 0.06 in proportion units. Hold confidence fixed at $\HLConfidence{95\%}$, using $z^*=z_{.95}$. 
\QuestionsTasks[sec:worked-solution-examples]{{6}/{sol:t3-prop}/{Find the minimum integer sample size. State the unknown, target, fixed quantities, units and model; rearrange the inequality, round correctly, check adjacent integers and interpret the result.}}
\subsubsection*{Worked solution (C6)}\label{sol:t3-prop}
\textbf{Design.} The unknown is $\HLObservationCount{n}$, the number of usable household responses. The target is total interval width at most 6 percentage points, namely 0.06 in proportion units with 95\% confidence. The planning share remains fixed; $p$ is unknown and $\hat p$ is the future estimate. Use the normal approximation for random binary outcomes; expected-count checks appear below. Confidence also remains fixed. 
\[
2\HLCriticalValue{z^*}\sqrt{\frac{\HLModelControl{0.50}(0.50)}{n}}\le \HLIntervalWidth{0.06}\quad\Longrightarrow\quad\sqrt n\ge\frac{2\HLCriticalValue{z^*}\sqrt{\HLModelControl{0.50}(0.50)}}{0.06}.
\]
\[n\ge\left(\frac{2\HLCriticalValue{z^*}\sqrt{\HLModelControl{0.50}(0.50)}}{0.06}\right)^2\approx1067.071894637,\qquad \HLObservationCount{n_{\min}}=\left\lceil 1067.071894637\right\rceil=\HLObservationCount{1068}.\]
Round upward, because rounding down would violate the inequality.
\textbf{Adjacent-integer check} (proportion units):
\[ W(1068)\approx0.059973924\le \HLIntervalWidth{0.06},\qquad W(1067)\approx0.060002021>0.06.\]
The permitted half-width is $0.03$ proportion units; $0.06$ is the full width, not the margin. Expected outcome counts are $np_0=534$ and $n(1-p_0)=534$, both at least 10. The sampling fraction is $1068/100{,}000=1.068\%<10\%$. Because the target quantity decreases with $n$, the checks establish minimality. Plan for $\HLObservationCount{1068}$ usable household responses. This conservative variability bound applies under the normal-formula planning model, not as exact coverage or guaranteed realized accuracy. After data collection, check actual outcome counts and the interval conditions. 
\secInicio[prob:t3-prop][sec:worked-solution-examples]

\clearpage
\subsection{Type 4: minimum sample size to support a confidence target}\label{prob:t4}
Supply a target $\HLConfidence{C_0}\in(0,1)$ \emph{and} a fixed positive precision cap $\HLMarginOfError{E_0}$. Keep variability and that cap fixed. Confidence alone cannot determine $n$: one must also limit the interval's half-width (or total width).

Under the normal planning model, the maximum central confidence associated with fixed half-width $E_0$ is
\[
\HLConfidence{C(n)}=2\Phi\!\left(\frac{E_0\sqrt n}{\sqrt V}\right)-1.
\]
Indeed the interval reaches $E_0/SE$ standard errors on each side of its center. Since $\Phi$ is increasing, with $z_0=\Phi^{-1}((1+C_0)/2)$,
\[
C(n)\ge C_0\ \Longleftrightarrow\ \frac{E_0\sqrt n}{\sqrt V}\ge z_0
\ \Longleftrightarrow\ n\ge\frac{z_0^2V}{E_0^2},\qquad
n_{\min}=\left\lceil z_0^2V/E_0^2\right\rceil.
\]
This is the Type 2 inequality at the target critical value, viewed as the confidence a fixed precision cap can support. At the chosen $n$, use $C_0$ or any confidence up to $C(n)$ while respecting the cap. It does not permit arbitrarily high confidence at the same precision.
\SectionProblemsTableOfContents{
\SectionProblemLink{prob:t4-mean}{Mean: cargo records}\\
\SectionProblemLink{prob:t4-prop}{Proportion: pilot-based planning}}
\medskip
\clearpage
\subsubsection{Cargo planning: support at least 95 percent confidence}\label{prob:t4-mean}
A logistics agency estimates the unknown population mean cargo-clearance time. It will draw a simple random sample without replacement from $N=100{,}000$ eligible cargo records. Hold the known population spread $\HLPopulationSD{\sigma}=\HLPopulationSD{12\text{ min}}$ fixed. Assume the sample mean is suitably approximated by a normal distribution. The margin of error cap is 2 minutes. Require at least $\HLConfidence{95\%}$, using $z_0=z_{.95}$. Keep the positive precision cap fixed while finding the minimum sample size that supports the confidence target. 
\QuestionsTasks[sec:worked-solution-examples]{{6}/{sol:t4-mean}/{Find the minimum integer sample size. State the unknown, target, fixed quantities, units and model; rearrange the inequality, round correctly, check adjacent integers and interpret the result.}}
\subsubsection*{Worked solution (C6)}\label{sol:t4-mean}
\textbf{Design.} The unknown is $\HLObservationCount{n}$, the number of usable cargo records. The target is margin of error at most 2 minutes with at least 95\% confidence. Population variability remains fixed. Use the known-$\sigma$ normal model for the sample mean. The margin cap also remains fixed. 
\[C(n)\ge C_0\ \Longleftrightarrow\ E_0\sqrt n/\sqrt V\ge \HLCriticalValue{z_0}\ \Longleftrightarrow\ \HLCriticalValue{z_0}SE(n)\le E_0.\]

\[
\HLCriticalValue{z_0}\frac{\HLPopulationSD{12}}{\sqrt n}\le \HLMarginOfError{2}\quad\Longrightarrow\quad\sqrt n\ge\frac{\HLCriticalValue{z_0}\HLPopulationSD{12}}{2}.
\]
\[n\ge\left(\frac{\HLCriticalValue{z_0}\HLPopulationSD{12}}{2}\right)^2\approx138.292517545,\qquad \HLObservationCount{n_{\min}}=\left\lceil 138.292517545\right\rceil=\HLObservationCount{139}.\]
Round upward, because rounding down would violate the inequality.
\textbf{Check confidence at fixed half-width.} With $V=144$,
\[
C(139)=2\Phi(2\sqrt{139}/\sqrt{144})-1\approx 0.950582408\ge 0.95,
\]
\[C(138)\approx 0.949757081<0.95.\]
Equivalently, at the target confidence the margins are $1.994903705\le \HLMarginOfError{2}$ and $2.002118570>2$ minutes, respectively. The sampling fraction is $139/100{,}000=0.139\%<10\%$. Because C(n) increases with $n$, the checks establish minimality. Plan for $\HLObservationCount{139}$ usable cargo records. 
\secInicio[prob:t4-mean][sec:worked-solution-examples]
\clearpage
\subsubsection{Subsidy planning: support at least 90 percent confidence}\label{prob:t4-prop}
An energy office estimates the unknown population share $p$ of households eligible for a bill subsidy. It will draw a simple random sample without replacement from $N=100{,}000$ households. A pilot suggests $\HLModelControl{p_0}=\HLModelControl{0.30}$; hold this planning value fixed. The margin of error cap is 3 percentage points, namely 0.03 in proportion units. Require at least $\HLConfidence{90\%}$, using $z_0=z_{.90}$. Keep the positive precision cap fixed while finding the minimum sample size that supports the confidence target. 
\QuestionsTasks[sec:worked-solution-examples]{{6}/{sol:t4-prop}/{Find the minimum integer sample size. State the unknown, target, fixed quantities, units and model; rearrange the inequality, round correctly, check adjacent integers and interpret the result.}}
\subsubsection*{Worked solution (C6)}\label{sol:t4-prop}
\textbf{Design.} The unknown is $\HLObservationCount{n}$, the number of usable household responses. The target is margin of error at most 3 percentage points, namely 0.03 in proportion units with at least 90\% confidence. The planning share remains fixed; $p$ is unknown and $\hat p$ is the future estimate. Use the normal approximation for random binary outcomes; expected-count checks appear below. The margin cap also remains fixed. 
\[C(n)\ge C_0\ \Longleftrightarrow\ E_0\sqrt n/\sqrt V\ge \HLCriticalValue{z_0}\ \Longleftrightarrow\ \HLCriticalValue{z_0}SE(n)\le E_0.\]

\[
\HLCriticalValue{z_0}\sqrt{\frac{\HLModelControl{0.30}(0.70)}{n}}\le \HLMarginOfError{0.03}\quad\Longrightarrow\quad\sqrt n\ge\frac{\HLCriticalValue{z_0}\sqrt{\HLModelControl{0.30}(0.70)}}{0.03}.
\]
\[n\ge\left(\frac{\HLCriticalValue{z_0}\sqrt{\HLModelControl{0.30}(0.70)}}{0.03}\right)^2\approx631.293472622,\qquad \HLObservationCount{n_{\min}}=\left\lceil 631.293472622\right\rceil=\HLObservationCount{632}.\]
Round upward, because rounding down would violate the inequality.
\textbf{Check confidence at fixed half-width.} With $V=0.21$,
\[
C(632)=2\Phi(0.03\sqrt{632}/\sqrt{0.21})-1\approx 0.900189663\ge 0.9,
\]
\[C(631)\approx 0.899921103<0.9.\]
Equivalently, at the target confidence the margins are $0.029983226\le \HLMarginOfError{0.03}$ and $0.030006976>0.03$ proportion units, respectively. Expected outcome counts are $np_0=189.6$ and $n(1-p_0)=442.4$, both at least 10. The sampling fraction is $632/100{,}000=0.632\%<10\%$. Because C(n) increases with $n$, the checks establish minimality. Plan for $\HLObservationCount{632}$ usable household responses. This pilot-based result is conditional on the assumed planning proportion. After data collection, check actual outcome counts and the interval conditions. 
\secInicio[prob:t4-prop][sec:worked-solution-examples]

\clearpage
\section{Advanced Problems}\label{sec:advanced}
\SectionProblemsTableOfContents{
\SectionProblemLink{prob:advanced-scale}{Compare precision and confidence targets}\\
\SectionProblemLink{prob:advanced-design}{Coordinate precision and confidence}\\
\SectionProblemLink{prob:advanced-planning}{Compare planning assumptions and justify conclusions}}
\medskip
\subsection{Compare precision and confidence targets}\label{prob:advanced-scale}
\QuestionsTasks[sec:advanced]{
{6}/{sol:advanced-scale}/{With variability fixed, explain how halving an SE cap changes the unrounded sample-size requirement. Repeat for a margin cap and a total-width cap, keeping confidence fixed. Does the integer answer always quadruple?},
{6}/{sol:advanced-confidence}/{For mean cargo-clearance time, hold known $\sigma=12$ minutes and $E_0=3$ minutes fixed. Compare 90\% and 95\% confidence. Use a random sample from $N=100{,}000$ records and a suitable normal model. Explain the direction of change and verify each minimum.}}
\subsubsection*{Solution: changing precision}\label{sol:advanced-scale}
For an SE cap, $b=V/S^2$; replacing $S$ by $S/2$ gives $4b$. The same reasoning applies to $b=(z^*)^2V/E_0^2$ and $b=4(z^*)^2V/W_0^2$ when the corresponding cap is halved and confidence stays fixed. Therefore the \emph{unrounded} requirement quadruples in all three cases.

The integer requirement changes from $\lceil b\rceil$ to $\lceil4b\rceil$, which need not equal $4\lceil b\rceil$. For example, the cargo margin problem has $b=61.463341131$: halving the margin changes the bound to $245.853364524$ and the minimum from 62 to 246, whereas $4(62)=248$. Distinguish the inverse-square relationship from the ceiling step.
\secInicio[prob:advanced-scale][sec:advanced]
\subsubsection*{Solution: increasing confidence at fixed precision}\label{sol:advanced-confidence}
The unknown is the usable-record count $n$. Keep $\sigma$ and the 3-minute half-width cap fixed. Under the known-$\sigma$ normal model,
\[
z^*12/\sqrt n\le3\ \Longleftrightarrow\ \sqrt n\ge4z^*
\ \Longleftrightarrow\ n\ge16(z^*)^2.
\]
Using the stored normal quantiles from Section 2 gives:
\begin{center}\begin{tabular}{@{}llll@{}}\toprule
Confidence & Unrounded bound & Ceiling & Adjacent margins (min)\\\midrule
90\% & 43.288695265 & 44 & $E(44)\approx2.975652$; $E(43)\approx3.010054$\\
95\% & 61.463341131 & 62 & $E(62)\approx2.98699$; $E(61)\approx3.01137$\\\bottomrule
\end{tabular}\end{center}
Each chosen integer meets the 3-minute cap and one fewer fails. Both samples are below 10\% of the population. Require 44 or 62 usable records respectively. Higher confidence has a larger $z^*$, so at fixed precision it requires more observations; the unrounded requirement is multiplied by $(z_{.95}/z_{.90})^2\approx1.41985$.
\secInicio[prob:advanced-scale][sec:advanced]
\clearpage
\subsection{Coordinate precision and confidence}\label{prob:advanced-design}
\QuestionsTasks[sec:advanced]{{6}/{sol:advanced-design}/{A cargo study has known $\sigma=12$ minutes, requires $SE\le1.5$ minutes and at least 95\% confidence with a fixed margin cap of 2 minutes. Under the same normal model, find the minimum $n$ satisfying both constraints. Use $N=100{,}000$.}}
\subsubsection*{Solution: two simultaneous constraints}\label{sol:advanced-design}
The unknown is the usable-record count. Keep variability, the SE cap and the positive half-width cap fixed. The known-$\sigma$ normal model gives
\[
12/\sqrt n\le1.5\ \Longleftrightarrow\ n\ge64,\qquad
z_{.95}12/\sqrt n\le2\ \Longleftrightarrow\ n\ge138.292517545.
\]
Take the ceiling of the larger bound: $n_{\min}=139$. At 139, $SE\approx1.017827$ minutes and $E_{.95}\approx1.994904$ minutes, so both caps hold. At 138, $E_{.95}\approx2.002119>2$ minutes; therefore no smaller sample works. The sampling fraction is $0.139\%$. Require 139 usable cargo records; the confidence calculation in Type 4 confirms the fixed-cap target as well.
\secInicio[prob:advanced-design][sec:advanced]
\clearpage
\subsection{Compare planning assumptions and justify conclusions}\label{prob:advanced-planning}
\QuestionsTasks[sec:advanced]{
{6}/{sol:advanced-pilot}/{For a random energy-eligibility survey from $N=100{,}000$ households, fix 95\% confidence and a margin cap of 4 percentage points. Compare the pilot choice $p_0=0.20$ with conservative $p_0=0.50$. Explain what each plan supports and what must be checked later.},
{6}/{sol:advanced-missing}/{An analyst requests at least 95\% confidence but specifies no SE, margin or width cap. Explain why this does not determine a minimum sample size for precision planning.}}
\subsubsection*{Solution: pilot versus conservative planning}\label{sol:advanced-pilot}
The unknown is the number of completed responses. Keep confidence and $E_0=0.04$ fixed; compare two assumed variability inputs in the binary normal planning model. Starting with
\[
z_{.95}\sqrt{p_0(1-p_0)/n}\le0.04
\ \Longleftrightarrow\ \sqrt n\ge\frac{z_{.95}\sqrt{p_0(1-p_0)}}{0.04}
\ \Longleftrightarrow\ n\ge\frac{z_{.95}^2p_0(1-p_0)}{0.04^2},
\]
we obtain:
\begin{center}\begin{tabular}{@{}llll@{}}\toprule
$p_0$ & Bound & Minimum & Adjacent planning margins\\\midrule
0.20 & 384.145882069 & 385 & $E(385)\approx0.039955606$; $E(384)\approx0.040007597$\\
0.50 & 600.227940733 & 601 & $E(601)\approx0.039974299$; $E(600)\approx0.040007597$\\\bottomrule
\end{tabular}\end{center}
Each ceiling meets 0.04 and one fewer fails under its planning assumption. Expected counts are 77 and 308 for the pilot plan, and 300.5 and 300.5 for the conservative plan. Sampling fractions are $0.385\%$ and $0.601\%$, both below 10\%.

Plan for 385 completed responses conditionally on $p_0=0.20$, or 601 using the maximum product $p_0(1-p_0)=0.25$. The conservative plan needs 216 extra responses; its unrounded requirement is $0.25/0.16=1.5625$ times as large. This is a bound within the normal-formula planning model, not a guarantee of exact coverage or realized accuracy. Check actual outcome counts and interval conditions after collection.
\secInicio[prob:advanced-planning][sec:advanced]
\subsubsection*{Solution: a missing precision constraint}\label{sol:advanced-missing}
Confidence specifies the multiplier, not an allowed width. Within the normal model, many valid sample sizes can produce a 95\% interval by allowing $E(n)=z_{.95}\sqrt{V/n}$ to adjust. A larger interval can accommodate the confidence target with fewer observations. Consequently confidence alone gives no precision-based minimum. Supply a positive fixed cap $E_0$ (or $W_0$), together with variability, before solving $C(n)\ge C_0$. Model-validity requirements are separate from a specified precision target.
\secInicio[prob:advanced-planning][sec:advanced]
\end{document}

